\documentclass[12pt]{article}
\usepackage[utf8]{inputenc}
\usepackage[T1]{fontenc}
\usepackage[margin=1in]{geometry}
\usepackage{setspace}
\usepackage{xcolor}
\usepackage[colorlinks=true, urlcolor=blue!60!black, linkcolor=blue!60!black]{hyperref}
\setlength{\parindent}{0em}
\setlength{\parskip}{0.6em}
\renewcommand{\baselinestretch}{1.0}
\begin{document}
\begin{center}
{\LARGE \textit{Learning Resources, Inc. v. Trump}, 607 U.S. \_\_\_ (2026)\par}
\vspace{0.5em}
{\large Case Brief --- Two-Page Edition\par}
\end{center}

\section{Memory Jogger}
IEEPA's power to ``regulate \ldots{} importation'' does not include the power to impose tariffs; the Justices split over the Major Questions Doctrine (\href{https://www.supremecourt.gov/opinions/25pdf/24-1287_4gcj.pdf}{Slip Op.}).

\section{Detailed Case Facts}
In 2025 President Trump invoked IEEPA to impose ``drug trafficking'' tariffs on Canada, Mexico, and China and worldwide ``reciprocal'' tariffs of at least 10\% (reaching 145\% on most Chinese goods), collecting more than \$200 billion. Two family-owned toy and education companies, and in a parallel case small businesses and twelve States, sued.

\section{Procedural History}
The D.C. district court enjoined the tariffs; the Court of International Trade and the en banc Federal Circuit held them unauthorized. The Supreme Court consolidated the cases and decided them on February 20, 2026.

\section{Judicial Votes}
6--3. Roberts wrote for Sotomayor, Kagan, Gorsuch, Barrett, and Jackson on the statutory holding, but only Gorsuch and Barrett joined his major-questions analysis; Gorsuch, Barrett, Kagan, and Jackson concurred; Kavanaugh (with Thomas and Alito) and Thomas dissented.

\section{Holding}
IEEPA does not authorize tariffs. The tariff power belongs to Congress, and ``regulate'' does not include the distinct power to raise revenue. The \textit{Learning Resources} judgment was vacated because the Court of International Trade's jurisdiction is exclusive.

\section{Analysis of Opinions}
The majority reasoned that IEEPA never mentions tariffs or duties, that ``regulate'' does not ordinarily mean tax, that a taxing reading would make export regulation unconstitutional, and that IEEPA's verbs describe sanctions. The plurality added the Major Questions Doctrine. Gorsuch defended it as a ``pro-Congress'' clear-statement rule; Barrett called it ordinary contextual textualism; Kagan found ordinary construction enough. Kavanaugh's dissent argued that tariffs are a classic means to regulate importation and that foreign affairs warrants deference; Thomas would limit nondelegation to core private rights.

\section{Examples: Future Applications}
\textbf{Same-Side:} a carbon border tax on imported cars and steel imposed under IEEPA would fall. \textbf{Opposite-Side:} IEEPA asset freezes and import embargoes against an adversary remain valid. \textbf{Fence-Sitter:} a cost-recovery ``licensing fee'' on imported AI models, or tariffs during an authorized war.

\section{Critique}
Abbe Gluck described a ``canon war'' between Gorsuch and Barrett over what the Major Questions Doctrine is. Ronald Levin argued that the decision left the doctrine ``more tenuous,'' since only three Justices joined it and Kagan showed it unnecessary.

\section{Key Quotations}
\begin{quote}
``Based on two words separated by 16 others in Section 1702(a)(1)(B) of IEEPA---`regulate' and `importation'---the President asserts the independent power to impose tariffs on imports from any country, of any product, at any rate, for any amount of time. Those words cannot bear such weight.''
\end{quote}
\noindent --- Chief Justice Roberts, Opinion of the Court, slip op. at 5.

\end{document}
